\documentclass[tikz,border=8pt]{standalone}
\usepackage{tikz}
\usepackage{amsmath}
\usepackage{amssymb}
\usepackage{pgfplots}
\pgfplotsset{compat=1.18}
\usepackage{ctex}
\usetikzlibrary{calc,positioning,arrows.meta}

% LC 167 Two Sum II — 对撞双指针
% 升序数组 [2,7,11,15]，target=9，3 个快照纵向排列

\begin{document}
\begin{tikzpicture}

\def\cellsize{0.8}   % 格子边长
\def\gap{2.8}        % 快照纵向间距

%% ── 辅助：画一个格子 ─────────────────────────────────────
% \drawcell{x}{y}{value}{fillcolor}
\newcommand{\drawcell}[4]{%
  \node[draw, minimum size=\cellsize cm, inner sep=0pt,
        fill=#4, font=\small\bfseries] at (#1,#2) {#3};
}

%% ── 标题 ─────────────────────────────────────────────────
\node[font=\large\bfseries] at (3.0, 0.8)
  {LC 167：对撞双指针（Two Sum II）};
\node[font=\small, gray] at (3.0, 0.3)
  {升序数组 [2, 7, 11, 15]，target = 9};

%% ═════════════════════════════════════════════════════════
%% 快照 0：初始，left=0，right=3，sum=2+15=17 > target
%% ═════════════════════════════════════════════════════════
\begin{scope}[yshift=-0.5cm]
  % 数组格子
  \drawcell{0}{0}{2}{red!15}     % index 0 — left
  \drawcell{0.9}{0}{7}{white}
  \drawcell{1.8}{0}{11}{white}
  \drawcell{2.7}{0}{15}{blue!15} % index 3 — right
  % 下标
  \foreach \i/\v in {0/0,1/1,2/2,3/3}{
    \node[font=\tiny, gray] at (\i*0.9, -0.55) {\v};
  }
  % left 箭头（红 ↑）
  \draw[->,>=Stealth,red!80!black,thick]
    (0, -0.9) -- (0, -0.5);
  \node[font=\scriptsize,red!80!black,below] at (0,-0.95) {left};
  % right 箭头（蓝 ↑）
  \draw[->,>=Stealth,blue!80!black,thick]
    (2.7, -0.9) -- (2.7, -0.5);
  \node[font=\scriptsize,blue!80!black,below] at (2.7,-0.95) {right};
  % 注释框
  \node[draw=gray!50,fill=white,thin,rounded corners=2pt,
        font=\scriptsize,inner sep=4pt,align=left]
    at (5.5, -0.3)
    {sum = 2+15 = \textbf{17} $>$ target=9\\
     $\Rightarrow$ right-- （缩小和）};
\end{scope}

%% ═════════════════════════════════════════════════════════
%% 快照 1：left=0，right=2，sum=2+11=13 > target
%% ═════════════════════════════════════════════════════════
\begin{scope}[yshift=-2.8cm]
  \drawcell{0}{0}{2}{red!15}
  \drawcell{0.9}{0}{7}{white}
  \drawcell{1.8}{0}{11}{blue!15}
  \drawcell{2.7}{0}{15}{gray!15}
  \foreach \i/\v in {0/0,1/1,2/2,3/3}{
    \node[font=\tiny, gray] at (\i*0.9, -0.55) {\v};
  }
  \draw[->,>=Stealth,red!80!black,thick]
    (0, -0.9) -- (0, -0.5);
  \node[font=\scriptsize,red!80!black,below] at (0,-0.95) {left};
  \draw[->,>=Stealth,blue!80!black,thick]
    (1.8, -0.9) -- (1.8, -0.5);
  \node[font=\scriptsize,blue!80!black,below] at (1.8,-0.95) {right};
  \node[draw=gray!50,fill=white,thin,rounded corners=2pt,
        font=\scriptsize,inner sep=4pt,align=left]
    at (5.5, -0.3)
    {sum = 2+11 = \textbf{13} $>$ target=9\\
     $\Rightarrow$ right-- （继续缩小）};
\end{scope}

%% ═════════════════════════════════════════════════════════
%% 快照 2：left=0，right=1，sum=2+7=9 == target → 找到！
%% ═════════════════════════════════════════════════════════
\begin{scope}[yshift=-5.6cm]
  \drawcell{0}{0}{2}{red!15}
  \drawcell{0.9}{0}{7}{blue!15}
  \drawcell{1.8}{0}{11}{gray!15}
  \drawcell{2.7}{0}{15}{gray!15}
  \foreach \i/\v in {0/0,1/1,2/2,3/3}{
    \node[font=\tiny, gray] at (\i*0.9, -0.55) {\v};
  }
  \draw[->,>=Stealth,red!80!black,thick]
    (0, -0.9) -- (0, -0.5);
  \node[font=\scriptsize,red!80!black,below] at (0,-0.95) {left};
  \draw[->,>=Stealth,blue!80!black,thick]
    (0.9, -0.9) -- (0.9, -0.5);
  \node[font=\scriptsize,blue!80!black,below] at (0.9,-0.95) {right};
  % 绿色框标记答案
  \draw[green!60!black,thick,rounded corners=2pt]
    (-0.42, -0.42) rectangle (1.32, 0.42);
  \node[font=\tiny,green!60!black] at (0.45, 0.6) {答案};
  \node[draw=gray!50,fill=white,thin,rounded corners=2pt,
        font=\scriptsize,inner sep=4pt,align=left]
    at (5.5, -0.3)
    {sum = 2+7 = \textbf{9} $=$ target\\
     $\Rightarrow$ 返回 [\textbf{1}, \textbf{2}]（1-indexed）};
\end{scope}

%% ── 图例 ─────────────────────────────────────────────────
\begin{scope}[yshift=-8.1cm]
  \node[draw,minimum size=0.5cm,fill=red!15,inner sep=0pt]  at (0.0,0) {};
  \node[font=\scriptsize,right] at (0.35,0) {left 指针（红）};
  \node[draw,minimum size=0.5cm,fill=blue!15,inner sep=0pt] at (3.2,0) {};
  \node[font=\scriptsize,right] at (3.55,0) {right 指针（蓝）};
  \node[draw,minimum size=0.5cm,fill=gray!15,inner sep=0pt] at (6.2,0) {};
  \node[font=\scriptsize,right] at (6.55,0) {已排除};
\end{scope}

\end{tikzpicture}
\end{document}
